% Options for packages loaded elsewhere
% Options for packages loaded elsewhere
\PassOptionsToPackage{unicode}{hyperref}
\PassOptionsToPackage{hyphens}{url}
%
\documentclass[
  ignorenonframetext,
  aspectratio=169,
]{beamer}
\newif\ifbibliography
\usepackage{pgfpages}
\setbeamertemplate{caption}[numbered]
\setbeamertemplate{caption label separator}{: }
\setbeamercolor{caption name}{fg=normal text.fg}
\beamertemplatenavigationsymbolsempty
% remove section numbering
\setbeamertemplate{part page}{
  \centering
  \begin{beamercolorbox}[sep=16pt,center]{part title}
    \usebeamerfont{part title}\insertpart\par
  \end{beamercolorbox}
}
\setbeamertemplate{section page}{
  \centering
  \begin{beamercolorbox}[sep=12pt,center]{section title}
    \usebeamerfont{section title}\insertsection\par
  \end{beamercolorbox}
}
\setbeamertemplate{subsection page}{
  \centering
  \begin{beamercolorbox}[sep=8pt,center]{subsection title}
    \usebeamerfont{subsection title}\insertsubsection\par
  \end{beamercolorbox}
}
% Prevent slide breaks in the middle of a paragraph
\widowpenalties 1 10000
\raggedbottom
\usepackage{iftex}
\ifPDFTeX
  \usepackage[T1]{fontenc}
  \usepackage[utf8]{inputenc}
  \usepackage{textcomp} % provide euro and other symbols
\else % if luatex or xetex
  \usepackage{unicode-math} % this also loads fontspec
  \defaultfontfeatures{Scale=MatchLowercase}
  \defaultfontfeatures[\rmfamily]{Ligatures=TeX,Scale=1}
\fi
\usepackage{lmodern}

\usetheme[]{metropolis}
\usecolortheme[]{default}
\usefonttheme[]{default}
\ifPDFTeX\else
  % xetex/luatex font selection
\fi
% Use upquote if available, for straight quotes in verbatim environments
\IfFileExists{upquote.sty}{\usepackage{upquote}}{}
\IfFileExists{microtype.sty}{% use microtype if available
  \usepackage[]{microtype}
  \UseMicrotypeSet[protrusion]{basicmath} % disable protrusion for tt fonts
}{}
\makeatletter
\@ifundefined{KOMAClassName}{% if non-KOMA class
  \IfFileExists{parskip.sty}{%
    \usepackage{parskip}
  }{% else
    \setlength{\parindent}{0pt}
    \setlength{\parskip}{6pt plus 2pt minus 1pt}}
}{% if KOMA class
  \KOMAoptions{parskip=half}}
\makeatother


\usepackage{longtable,booktabs,array}
\usepackage{calc} % for calculating minipage widths
\usepackage{caption}
% Make caption package work with longtable
\makeatletter
\def\fnum@table{\tablename~\thetable}
\makeatother
\usepackage{graphicx}
\makeatletter
\newsavebox\pandoc@box
\newcommand*\pandocbounded[1]{% scales image to fit in text height/width
  \sbox\pandoc@box{#1}%
  \Gscale@div\@tempa{\textheight}{\dimexpr\ht\pandoc@box+\dp\pandoc@box\relax}%
  \Gscale@div\@tempb{\linewidth}{\wd\pandoc@box}%
  \ifdim\@tempb\p@<\@tempa\p@\let\@tempa\@tempb\fi% select the smaller of both
  \ifdim\@tempa\p@<\p@\scalebox{\@tempa}{\usebox\pandoc@box}%
  \else\usebox{\pandoc@box}%
  \fi%
}
% Set default figure placement to htbp
\def\fps@figure{htbp}
\makeatother





\setlength{\emergencystretch}{3em} % prevent overfull lines

\providecommand{\tightlist}{%
  \setlength{\itemsep}{0pt}\setlength{\parskip}{0pt}}



 


\makeatletter
\@ifpackageloaded{caption}{}{\usepackage{caption}}
\AtBeginDocument{%
\ifdefined\contentsname
  \renewcommand*\contentsname{Table of contents}
\else
  \newcommand\contentsname{Table of contents}
\fi
\ifdefined\listfigurename
  \renewcommand*\listfigurename{List of Figures}
\else
  \newcommand\listfigurename{List of Figures}
\fi
\ifdefined\listtablename
  \renewcommand*\listtablename{List of Tables}
\else
  \newcommand\listtablename{List of Tables}
\fi
\ifdefined\figurename
  \renewcommand*\figurename{Figure}
\else
  \newcommand\figurename{Figure}
\fi
\ifdefined\tablename
  \renewcommand*\tablename{Table}
\else
  \newcommand\tablename{Table}
\fi
}
\@ifpackageloaded{float}{}{\usepackage{float}}
\floatstyle{ruled}
\@ifundefined{c@chapter}{\newfloat{codelisting}{h}{lop}}{\newfloat{codelisting}{h}{lop}[chapter]}
\floatname{codelisting}{Listing}
\newcommand*\listoflistings{\listof{codelisting}{List of Listings}}
\makeatother
\makeatletter
\makeatother
\makeatletter
\@ifpackageloaded{caption}{}{\usepackage{caption}}
\@ifpackageloaded{subcaption}{}{\usepackage{subcaption}}
\makeatother

\usepackage{bookmark}
\IfFileExists{xurl.sty}{\usepackage{xurl}}{} % add URL line breaks if available
\urlstyle{same}
\hypersetup{
  hidelinks,
  pdfcreator={LaTeX via pandoc}}


\author{}
\date{}

\begin{document}


\section{\texorpdfstring{\textbf{Presentation Outline: A Cohesive System
for Engineering
Leverage}}{Presentation Outline: A Cohesive System for Engineering Leverage}}\label{presentation-outline-a-cohesive-system-for-engineering-leverage}

\begin{frame}{\textbf{Title Slide}}
\phantomsection\label{title-slide}
\begin{itemize}
\tightlist
\item
  \textbf{Automating the Mundane, Building the Monumental: A
  Practitioner's System for AI-Driven Engineering}\\
\item
  Presented by: Mo Battah\\
\item
  Audience: {[}Target Engineering Team/Conference{]}\\
\item
  Date: 08/20/2025\\
\item
  Link to Open Source Repo: {[}Link to your Git Repository{]}
\end{itemize}
\end{frame}

\begin{frame}{\textbf{Part 1: The Philosophy (The ``Why'')}}
\phantomsection\label{part-1-the-philosophy-the-why}
\emph{(Goal: Establish immediate resonance with fellow engineers by
framing the problem as engineering toil vs.~creative work. Present a
philosophy of work that every engineer will covet.)}

\begin{block}{\textbf{A. Introduction: Speaking Peer-to-Peer}}
\phantomsection\label{a.-introduction-speaking-peer-to-peer}
\begin{itemize}
\item
  \textbf{My Background:} Recently VP, DevOps Engineering at
  Actabl/Alpine Investors. Since leaving 1.5 months ago, I've advised on
  \textgreater\$10M in transactions and am currently helping founders in
  AI agents, cybersecurity, and deep tech.
\item
  \textbf{Portfolio Context for Speaking Notes:}

  \begin{itemize}
  \tightlist
  \item
    \textbf{ProtectAI} (→ Palo Alto Networks): AI security, protecting
    ML models from attacks
  \item
    \textbf{Skydio}: Autonomous drones revolutionizing power line
    inspection, bringing down electricity costs
  \item
    \textbf{Dreadnode}: Red teaming against AI systems, offensive LLM
    security for government and enterprise
  \item
    \textbf{Resourcely} (→ Cursor): Infrastructure as code platform
  \item
    \textbf{Kodiak Robotics}: Autonomous trucking already operating in
    Permian Basin, moving sand for fracking operations
  \item
    \textbf{Additional notable portfolio}: Sublime Security (trusted by
    Netflix, Spotify), Nucleus Security (American Airlines, Mastercard),
    SpecterOps (market-leading cybersecurity), and many others
  \end{itemize}
\item
  \textbf{AI as Deflationary Force Context}: Skydio and Kodiak are
  perfect examples of AI dramatically reducing costs in critical
  infrastructure. Skydio is making power line inspection vastly cheaper,
  potentially reducing electricity costs. Kodiak's autonomous trucks in
  the Permian Basin are already reducing transportation costs for oil
  operations, potentially affecting gas and oil prices through
  efficiency gains.
\item
  \textbf{Refined Narrative:} ``AI isn't about replacing us; it's about
  automating the mundane so we can focus on the monumental. I'm not here
  to show you theoretical futures---I'm showing you systems I use daily
  to close real deals and help portfolio companies scale.''
\end{itemize}
\end{block}

\begin{block}{\textbf{B. Core Principle: Augmentation, Not Abdication}}
\phantomsection\label{b.-core-principle-augmentation-not-abdication}
\begin{itemize}
\tightlist
\item
  \textbf{Core Message:} This is not about blindly trusting AI or
  outsourcing our thinking. It's about a fundamental shift in our role,
  moving from a ``doer'' to an ``architect.'' It's a new, more leveraged
  form of engineering leadership where we define the problem, design the
  solution, and act as the ultimate verifiers of quality.\\
\item
  \textbf{The Key Axiom:} ``If a task can be fully automated to a
  satisfactory degree, our job is to find the next, more valuable task
  up the chain. We use AI to augment our judgment, not replace it. The
  goal is to operate at a higher level of abstraction.''
\end{itemize}
\end{block}
\end{frame}

\begin{frame}{\textbf{Part 2: The Practitioner's Stack (The ``How'')}}
\phantomsection\label{part-2-the-practitioners-stack-the-how}
\emph{(Goal: Showcase your personal toolkit not as a random collection
of apps, but as a deliberately architected stack for achieving maximum
leverage.)}

\begin{block}{\textbf{A. My Stack for Leverage: A 4-Stage Context
Engineering Pipeline}}
\phantomsection\label{a.-my-stack-for-leverage-a-4-stage-context-engineering-pipeline}
\emph{(Core Message: Context is gold. Without good context engineering,
prompt engineering is irrelevant. But here's the deeper truth---it's not
just prompt engineering. All AI outputs are completely irrelevant
without good context. Context poisoning is the real threat. This isn't a
random collection of tools; it's a deliberate, four-stage pipeline for
manufacturing perfect, context-rich prompts. High-quality output
requires high-quality input, and this is the system I use to create
it.)}

\textbf{Key Quotes to Reference:} - \emph{``Context is everything, yet
context is what's missing.''} - This paradox captures the challenge
perfectly. Context determines everything about AI quality, yet it's
consistently the weakest link in most people's AI workflows. -
\emph{``The whole world is an intelligence game.''} - A friend's insight
that became crystal clear once I understood context engineering. In
business, in technology, in relationships---it's all about having better
information and context than the other party. AI amplifies this reality
exponentially.

\textbf{Speaking Note:} ``These quotes didn't make sense to me until I
got deep into context engineering. Now I realize that context
engineering isn't just a technical skill---it's the fundamental skill
for the AI age. Bad context doesn't just give you bad outputs; it
poisons your entire decision-making process.''

\begin{itemize}
\tightlist
\item
  \textbf{Stage 1: Context Inception (Wispr Flow)}

  \begin{itemize}
  \tightlist
  \item
    \textbf{Pitch:} ``The entire process starts with high-bandwidth
    thought capture. Keyboards are a bottleneck. Wispr Flow is my
    primary compute interface. Wispr allows you to talk/dictate/speak
    into any textbox whether on your computer or phone. Already, this is
    a significant speed advancement, People will naturally output more
    text if the input mechanism is speech instead of typing on a
    keyboard. Now, where this goes from a handy daily tool to the base
    of context engineering is that you can use Wispr Flow for turning
    unstructured verbal brainstorming into structured, AI-ready text.
    Your thoughts, you just speak, verbalize in real-time, transcribed
    become a digital asset---the raw material for the entire AI
    workflow.''\\
  \item
    \textbf{Application:} ``This is foundational. I use it for
    everything from drafting emails to capturing real-time thoughts
    while walking through a house, or more importantly, talking through
    a complex systems design and architecture problem at work. The
    output is a repository of authentic, high-fidelity text that becomes
    our initial, unvarnished context. Now, you might ask: Hey Mo, you're
    telling me I speak and this technology will listen to me and
    transcribe my thoughts, and maybe I'll even use AI to help structure
    those thoughts? Okay, I get you. Well Mo, where do I store all these
    documents? Do I store them in Google Docs or Word? Let's talk about
    context storage, and management.''
  \item
    \textbf{Volume:} \textasciitilde100K words/month of voice dictation
  \item
    \textbf{Key Insight:} Your thoughts become digital assets that can
    be processed and leveraged by AI
  \end{itemize}
\item
  \textbf{Stage 2: Context Storage \& Management (Git Markdown
  Repository as Second Brain)}

  \begin{itemize}
  \tightlist
  \item
    \textbf{Pitch:} ``This is where we ensure deterministic, reliable
    context instead of relying on probabilistic MCP server calls. I
    download all emails with individual partners, clients, or business
    contacts into folders, then use Claude Code or Gemini CLI to flatten
    all the emails into markdown files that get stored in my Git
    repository.''\\
  \item
    \textbf{Why GitHub over Notion:} ``I used to put everything in
    Notion, but the Notion MCP server is too finicky. Whenever you're
    using an MCP server, you're relying on probabilistic outcomes.
    Instead of asking AI to always review the latest emails with someone
    through MCP tool calls---which you sometimes don't know if it messes
    up or not---I just have a file of all the emails. That way it's
    deterministic, the data is there, and all AI has to do is help me
    with the task I'm asking, which is usually a question about the
    context.''
  \item
    \textbf{Application:} Email histories become searchable, reliable
    context files that AI can always access
  \item
    \textbf{Technical Details:} Download email folders → AI flattens to
    markdown → commit to GitHub repo
  \item
    \textbf{Context Examples:} Project goals, business insights, client
    data, personal career strategy, email histories, call transcripts
  \end{itemize}
\item
  \textbf{Stage 3: Context Refinement \& Enrichment (Agentic Partners)}

  \begin{itemize}
  \tightlist
  \item
    \textbf{Pitch:} ``Raw text and stored context are just starting
    points. The next step is to refine and structure it. I use agentic
    partners like Claude Code and Gemini CLI as infinitely patient
    junior developers. They take the raw context from Wispr plus the
    stored context files, organize it, enrich it with research, and
    build it into more robust, analytical documents that are ready for
    complex queries.''\\
  \item
    \textbf{Application:} ``For that house tour, the agent transforms my
    rambling notes into a structured pros/cons list and researches
    neighborhood comps and school ratings. For a coding problem, it
    converts my verbal stream-of-consciousness into a preliminary bug
    report, complete with code snippets and a list of potential root
    causes. For business relationships, it can cross-reference my notes
    with the complete email history to provide strategic insights.''
  \end{itemize}
\item
  \textbf{Stage 4: Context Assembly \& Orchestration (Prompt Tower)}

  \begin{itemize}
  \tightlist
  \item
    \textbf{Pitch:} ``This is the final and most powerful step where we
    achieve insurmountable leverage. Prompt Tower is the ultimate
    context assembler. It takes all the disparate pieces of context
    we've created and refined---the Wispr transcripts, the
    agent-generated reports, the flattened email files, relevant code
    files, API documentation, meeting notes---and allows me to
    orchestrate the perfect, holistic prompt.''\\
  \item
    \textbf{Application:} ``For strategic call preparation---my most
    common use case---I select all past call transcripts, the flattened
    file of our email history, and a natural-language `status' file that
    captures the nuances of the relationship that a tool like Monday.com
    can't. My prompt is then simple: `Based on all this, what are my
    next steps? And critically, what potential dissatisfactions do I
    need to surface and address?' For the house purchase, I can select
    the refined reports for three properties, add a context file with my
    financial data, and another with my family's requirements. For a
    software engineering problem, I can select the bug report, the
    relevant codebase files, and transcripts from team calls discussing
    the issue. Prompt Tower assembles all of this into one massive,
    perfectly-formed prompt that gives the AI a complete, 360-degree
    view of the problem, enabling it to provide a deeply insightful and
    immediately useful solution. Typically, these prompts end up in
    Gemini because of their 1 million token context. My submissions are
    typically north of 100,000 tokens.''
  \item
    \textbf{Technical Features:} VSCode extension, file selection
    interface, custom prompt prefixes/suffixes, one-click clipboard copy
  \item
    \textbf{Economic Advantage:} Gemini's flat consumer fee allows
    unlimited iteration and prompt optimization
  \end{itemize}
\end{itemize}
\end{block}

\begin{block}{\textbf{B. The Demos: Primitives for Building Systems}}
\phantomsection\label{b.-the-demos-primitives-for-building-systems}
\emph{(Goal: Present your demos as case studies in applying engineering
principles to automate toil.)}

\begin{itemize}
\tightlist
\item
  \textbf{Demo 1: The Modular Resume System (A Case Study in
  ``Documentation as Code'')}

  \begin{itemize}
  \tightlist
  \item
    \textbf{The Pitch:} ``We all hate updating documentation because
    it's manual, repetitive, and prone to error. I applied a software
    engineering mindset to a classic business problem: the resume. This
    is a system built on modular, version-controlled components---our
    single source of truth---which are then assembled by templates and
    automated with build scripts. The result is publication-quality,
    ATS-optimized, and role-specific outputs generated on the fly.''\\
  \item
    \textbf{The Connection:} ``This is a powerful principle. Now,
    imagine applying it to your team's technical documentation, API
    specs, or customer-facing reports. We can stop manually writing docs
    and start generating them from the source, ensuring they are always
    accurate and up-to-date. This is how we solve the documentation
    problem for good.''\\
  \item
    \emph{(Showcase the resume-system/ demo. Run the build script.)}\\
  \end{itemize}
\item
  \textbf{Demo 2: AI-Assisted Technical Analysis (Automating the ``Last
  Mile'')}

  \begin{itemize}
  \tightlist
  \item
    \textbf{The Pitch:} ``This demonstrates the output of such a system,
    automating the final, painful step of presentation. It combines
    AI-driven content generation with the typographic precision of
    LaTeX. Notice the programmatically generated TikZ diagrams---no more
    struggling with Visio or Lucidchart. This entire professional-grade
    report was rendered from simple Markdown files with a single
    command: quarto render.''\\
  \item
    \textbf{The Connection:} ``This is how we automate the `last mile'
    of technical communication. It ensures that our hard work is
    presented with the professionalism it deserves, without hours of
    manual formatting. It allows us to focus on the substance of our
    analysis, knowing the system will handle the world-class
    presentation.''\\
  \item
    \emph{(Showcase the report-generation/ demo. Display the .qmd file
    next to the final PDF.)}\\
  \end{itemize}
\item
  \textbf{Demo 3: The AI-Powered Feedback Loop (Systematizing
  Self-Improvement)}

  \begin{itemize}
  \tightlist
  \item
    \textbf{The Pitch:} ``As engineers, we're great at getting feedback
    on our code, but often terrible at getting objective feedback on our
    communication. I treat this as another engineering problem. By
    recording and transcribing my professional calls (with consent), I
    create a dataset of my own performance. I then use AI to analyze it,
    asking critical questions: `Where are my communication gaps?', `How
    could I have handled this negotiation better?', `What is the
    sentiment of this call?' This has been one of the most high-leverage
    ways to improve myself. I will often ask AI for book
    recommendations---whether it is a new domain I'm getting into (like
    a certain sector of private equity), a specific communication skill
    that I'm interested in developing, or an interpersonal skill.''\\
  \item
    \textbf{The Connection:} ``This system provides a data-driven,
    objective feedback loop for the soft skills that determine career
    velocity. It helped me identify a tendency to be overly
    accommodative in negotiations and gave me the confidence to advocate
    for myself more effectively. It's a way to systematically debug our
    own professional performance, addressing the imposter syndrome that
    holds so many of us back.''\\
  \end{itemize}
\item
  \textbf{Demo 4: The AI Vocabulary Coach (Extending Expertise
  Infinitely)}

  \begin{itemize}
  \tightlist
  \item
    \textbf{The Pitch:} ``We often hit a plateau in our professional
    language. This system, inspired by legendary grammarians like Norman
    Lewis and Tom Heehler, acts as a personalized vocabulary coach. I
    feed it my call transcripts, and it analyzes my language to identify
    words just at the edge of my current vernacular, making my
    communication more precise and impactful. The prompts for this are
    open-sourced in my GitHub.''\\
  \item
    \textbf{The Connection:} ``This demonstrates a profound principle:
    AI can take a finite resource---like the knowledge in a beloved book
    or the expertise of a mentor who is no longer with us---and extend
    its value infinitely. It transforms a static resource into a
    dynamic, personalized learning engine that continuously refines our
    most critical tool: our language. This is how we systematize not
    just our work, but our own growth.''\\
  \end{itemize}
\item
  \textbf{Demo 5: The Automated Networking Strategist (Building an
  AI-Powered CRM)}

  \begin{itemize}
  \tightlist
  \item
    \textbf{The Pitch:} ``Networking is often reactive and based on
    memory. This system makes it proactive and strategic. When I get new
    connections on LinkedIn, I feed them to an agent. The agent
    researches each person and cross-references them against my personal
    context files---my career goals, business objectives, and areas of
    interest.''\\
  \item
    \textbf{The Connection:} ``The AI acts as a strategic filter. It
    reconciles the public data about a person with my private goals to
    determine if they are a relevant connection. If they are, it
    automatically adds them to my personal CRM---a simple Markdown
    database---with notes on why they're relevant. This transforms a
    passive activity into an automated, high-leverage networking
    strategy, ensuring I'm building a valuable network with almost zero
    manual effort.''\\
  \item
    \emph{(Show screenshot of the AI-generated CRM entry)}
  \end{itemize}
\end{itemize}
\end{block}
\end{frame}

\begin{frame}{\textbf{Part 3: The Paradigm Shift (The ``What If'')}}
\phantomsection\label{part-3-the-paradigm-shift-the-what-if}
\emph{(Goal: Connect your system directly to the daily work of the
audience and present a new, more powerful way of operating.)}

\begin{block}{\textbf{A. The New Paradigm: From Writing Code to
Directing Code Generation}}
\phantomsection\label{a.-the-new-paradigm-from-writing-code-to-directing-code-generation}
\begin{itemize}
\tightlist
\item
  \textbf{Core Message:} This is a fundamental shift in our role and
  where we provide value. The bottleneck is no longer typing speed or
  syntax knowledge; it's the quality of our thinking, our ability to
  design robust systems, and our rigor in verifying the output. The
  senior engineer's primary role becomes that of an architect and the
  final checkpoint for quality and correctness.\\
\item
  \textbf{The Anthropic PR is Your Proof Point:}

  \begin{itemize}
  \tightlist
  \item
    \textbf{Pitch:} ``This isn't a toy or a theoretical exercise. This
    is how high-performance teams are already operating. Anthropic
    successfully merged a 22,000-line, Claude-generated PR into their
    production RL systems. How? By shifting their entire review process
    away from painstaking line-by-line code review to a higher-level
    focus on behavioral verification and system-level testing.''\\
  \item
    \emph{(Show the ``Vibe Coding'' slides from index.qmd)}
  \end{itemize}
\end{itemize}
\end{block}

\begin{block}{\textbf{B. Principles for the New Paradigm}}
\phantomsection\label{b.-principles-for-the-new-paradigm}
\begin{itemize}
\tightlist
\item
  \textbf{1. Change the Abstraction:} Treat the AI like a
  hyper-productive team you manage. Your job is to write the spec, not
  the code. Focus on acceptance tests, stress tests, and verifiable
  system behaviors.\\
\item
  \textbf{2. Constrain the Blast Radius:} Start with ``leaf
  nodes''---isolated features, UI components, or scripts with no
  dependencies. Keep the core, foundational architecture human-owned and
  rigorously reviewed.\\
\item
  \textbf{3. Act as the Model's PM:} Spend 15-20 minutes creating a
  comprehensive brief. Your leverage comes from the quality of your
  direction---providing repository maps, context, constraints, and clear
  success criteria.\\
\item
  \textbf{4. Build for Verifiability:} Design for black-box testing. If
  you can't verify the output with a deterministic input/output contract
  and robust testing, you can't safely use the generated code.
\end{itemize}
\end{block}
\end{frame}

\begin{frame}{\textbf{Conclusion \& Discussion}}
\phantomsection\label{conclusion-discussion}
\begin{block}{\textbf{A. The Cohesive System}}
\phantomsection\label{a.-the-cohesive-system}
\begin{itemize}
\tightlist
\item
  \textbf{Closing Statement:} ``These aren't just isolated tricks or
  clever hacks. They are components of a cohesive system for
  professional practice, built on engineering principles. This is how I
  operate. It's a system designed to maximize the time we spend on
  interesting, high-impact problems and mercilessly eliminate the time
  we spend on everything else.''\\
\item
  \textbf{Subtle Positioning:} This subtly positions you as the expert
  they should turn to when they want to implement these ideas and build
  their own systems for leverage.
\end{itemize}
\end{block}

\begin{block}{\textbf{B. Answering Our Opening Questions}}
\phantomsection\label{b.-answering-our-opening-questions}
\begin{itemize}
\tightlist
\item
  \textbf{SDLC Change:} Engineers become architects and product managers
  for AI agents. The most critical skill is no longer writing code, but
  designing systems and verifying behavior. Review processes must shift
  from code-level to behavior-level.\\
\item
  \textbf{QA Evolution:} QA becomes even more critical, moving towards
  systems thinking, complex integration testing, and sophisticated
  black-box verification to validate business logic at scale.\\
\item
  \textbf{Team Ratios:} Senior architects who own the fundamental design
  and can maintain a coherent vision across an AI-augmented team become
  exponentially more valuable. The ratio of ``coders'' to ``architects''
  will likely shift.
\end{itemize}
\end{block}

\begin{block}{\textbf{C. Open Discussion}}
\phantomsection\label{c.-open-discussion}
\begin{itemize}
\tightlist
\item
  ``Now, back to you. What shifts do you see in your teams? What part of
  your process would you automate first? Let's talk.''
\end{itemize}
\end{block}

\begin{block}{\textbf{D. Final Slide: Resources}}
\phantomsection\label{d.-final-slide-resources}
\begin{itemize}
\tightlist
\item
  \textbf{Thank You}\\
\item
  Contact: {[}Your Email / LinkedIn{]}\\
\item
  \textbf{Open Source Repository:}

  \begin{itemize}
  \tightlist
  \item
    {[}Link to your Git Repository{]}\\
  \item
    Includes: This Presentation, Report Generation Demo, Modular Resume
    System, Prompt Library.
  \end{itemize}
\end{itemize}
\end{block}
\end{frame}




\end{document}
